\section{Failure Taxonomy}

\subsection{The Component-Based Taxonomy}
Rather than categorizing AI failures by subjective symptoms (e.g., ``model hallucinated''), FAULTLINE organizes failures by the \textbf{producing system component} where the defect originates:

\begin{center}
\resizebox{0.95\textwidth}{!}{%
\begin{tikzpicture}[
  rootbox/.style={draw=primary, fill=primary!10, rounded corners=3pt, font=\small\bfseries\sffamily, align=center, inner sep=5pt, text width=7.0cm},
  catbox/.style={draw=secondary, fill=secondary!10, rounded corners=3pt, font=\footnotesize\bfseries\sffamily, align=center, inner sep=4pt, text width=7.4cm},
  leafbox/.style={draw=accent, fill=accent!5, rounded corners=3pt, font=\scriptsize\sffamily, align=left, inner sep=4pt, text width=7.2cm},
  arrow/.style={-{Latex[length=2mm]}, line width=1pt, color=secondary},
  bus/.style={line width=1pt, color=secondary}
]
\node[rootbox] (root) at (0, 0) {FAULTLINE Failure Spectrum};

\node[catbox] (c1) at (-4.2, -1.2) {\textbf{Provider Boundary} (Transport and Interface)};
\node[catbox] (c2) at (4.2, -1.2) {\textbf{Agent Internal State} (Data Plane and Control)};

\node[leafbox, below=0.4cm of c1] (f1) {\textbf{F1: Structured Output Corruption}\\Truncated / invalid syntax at tool output boundary};
\node[leafbox, below=0.25cm of f1] (f2) {\textbf{F2: Virtual Latency / Timeout}\\Provider transport hang / queue saturation};
\node[leafbox, below=0.25cm of f2] (f4) {\textbf{F4: Provider Outage / HTTP Error}\\5xx error codes, 429 rate limits, service drop};

\node[leafbox, below=0.4cm of c2] (f3) {\textbf{F3: Schema Drift / Wrong Data}\\Syntactically valid payload citing invalid entity};
\node[leafbox, below=0.25cm of f3] (f5) {\textbf{F5: Context Corruption}\\Distractor injection, conflicting intermediate token};
\node[leafbox, below=0.25cm of f5] (f6) {\textbf{F6: Loop Exhaustion}\\Query cycle, pagination trap, step limit breach};

\draw[arrow] (root.south) -| (c1.north);
\draw[arrow] (root.south) -| (c2.north);

% Bus lines for c1
\draw[bus] (c1.south) -- ++(0,-0.2) -| ([xshift=-0.25cm]f1.north west) -- ([xshift=-0.25cm]f4.west);
\draw[arrow] ([xshift=-0.25cm]f1.west) -- (f1.west);
\draw[arrow] ([xshift=-0.25cm]f2.west) -- (f2.west);
\draw[arrow] ([xshift=-0.25cm]f4.west) -- (f4.west);

% Bus lines for c2
\draw[bus] (c2.south) -- ++(0,-0.2) -| ([xshift=-0.25cm]f3.north west) -- ([xshift=-0.25cm]f6.west);
\draw[arrow] ([xshift=-0.25cm]f3.west) -- (f3.west);
\draw[arrow] ([xshift=-0.25cm]f5.west) -- (f5.west);
\draw[arrow] ([xshift=-0.25cm]f6.west) -- (f6.west);
\end{tikzpicture}%
}
\vspace{0.2em}
\captionof{figure}{Taxonomy of AI Failure Modes by Producing Component.}
\label{fig:failure_taxonomy}
\end{center}

\begin{center}
\begingroup
\small
\setlength{\tabcolsep}{3pt}
\begin{tabularx}{\textwidth}{>{\raggedright\arraybackslash}p{0.18\textwidth}>{\raggedright\arraybackslash}p{0.18\textwidth}>{\raggedright\arraybackslash}p{0.18\textwidth}>{\raggedright\arraybackslash}p{0.20\textwidth}X}
\toprule
\textbf{Fault ID \& Component} & \textbf{Physical Trigger} & \textbf{Observable Trace} \newline \textbf{Symptom} & \textbf{User-Visible} \newline \textbf{Consequence} & \textbf{Detection \& Recovery} \\
\midrule
\textbf{F1: Structured Output} \newline \textit{Tool output boundary} & Truncated JSON, unescaped quotes, key omission. & Pydantic \code{ValidationError} at deserialization. & Unhandled crash or blank error message. & \textbf{Det:} Schema check ($R=1.0$). \newline \textbf{Rec:} M1 Schema repair-retry. \\
\addlinespace
\textbf{F2: Latency / Timeout} \newline \textit{Provider transport} & Simulated provider hang, packet drop, slow queue. & Span duration breaches configured latency threshold. & Client disconnect, spinning loader, SLA breach. & \textbf{Det:} Deadline monitor ($R=1.0$). \newline \textbf{Rec:} M2 Timeout + bounded retry. \\
\addlinespace
\textbf{F3: Schema Drift / Wrong Data} \newline \textit{Tool data plane} & Data format changes; field valid string but wrong entity. & Schema passes; downstream lookup fails to find entity. & Plausible hallucination citing invalid entity. & \textbf{Det:} Invariant check ($R=0.6$). \newline \textbf{Rec:} M6 Replanning / re-search. \\
\addlinespace
\textbf{F4: Provider Outage} \newline \textit{Provider infrastructure} & 500 Internal Error, 503 Overloaded, 429 Rate Limit. & Explicit error code returned by HTTP client wrapper. & Immediate query failure; cascade collapse. & \textbf{Det:} Status code check ($R=1.0$). \newline \textbf{Rec:} M3 Breaker + M4 Fallback. \\
\addlinespace
\textbf{F5: Context Corruption} \newline \textit{Agent memory / context} & Distractor injection, conflicting intermediate token. & Agent generates divergent tool queries unrelated to task. & Confident incorrect answer citing distractor. & \textbf{Det:} Semantic invariant ($R=0.5$). \newline \textbf{Rec:} Context reset / re-ground. \\
\addlinespace
\textbf{F6: Loop Exhaustion} \newline \textit{Agent control loop} & Query cycle (searching A $\to$ B $\to$ A), pagination trap. & Consecutive identical tool calls; step count reaches cap. & Silent timeout or partial ungrounded guess. & \textbf{Det:} Cycle counter ($R=1.0$). \newline \textbf{Rec:} M5 Structural step ceiling. \\
\bottomrule
\end{tabularx}
\endgroup
\end{center}

% ==============================================================================
\vspace{12pt}
\needspace{8\baselineskip}
\section{Fault Injection}

\subsection{Controlled Deterministic Injection}
In distributed systems, testing failure handling by waiting for production outages is unacceptable. Testing must be proactive and repeatable.

FAULTLINE injects failures deterministically at runtime:
\begin{center}
\resizebox{0.92\textwidth}{!}{%
\begin{tikzpicture}[
  node distance=0.6cm and 0.9cm,
  box/.style={draw=primary, fill=lightbg, rounded corners=3pt, font=\small\sffamily, align=center, inner sep=6pt},
  fault/.style={draw=accent, fill=accent!10, rounded corners=3pt, font=\small\bfseries\sffamily, align=center, inner sep=6pt},
  arrow/.style={-{Latex[length=2.5mm]}, line width=1pt, color=primary}
]
\node[box] (call) {Agent Emits\\Tool/Model Call};
\node[fault, right=of call] (inj) {Fault Injector\\Interception Proxy};
\node[box, right=of inj] (dest) {Target\\Component};
\node[fault, below=0.9cm of inj] (log) {Out-of-Band\\Ground Truth Log};

\draw[arrow] (call) -- (inj);
\draw[arrow] (inj) -- node[above, font=\scriptsize] {If healthy} (dest);
\draw[arrow, color=accent] (inj) -- node[right, font=\scriptsize] {Record injection} (log);
\draw[arrow, color=accent] (inj.south) -- ++(0,-0.35) -| node[below, pos=0.35, font=\scriptsize] {If faulted: inject F1--F6} (call.south);
\end{tikzpicture}%
}
\end{center}

\subsection{Independent Out-of-Band Ground Truth}
A foundational vulnerability in evaluation research is \textbf{ground-truth leakage}: if the detector inspects metadata injected into the data payload, the detector achieves an artificially perfect score that collapses in production.

In FAULTLINE, the injector logs the exact fault type, severity, target step, and timestamp to an \textbf{independent out-of-band manifest} (\code{injection\_truth.json}). The detector operates with zero knowledge of this file, inspecting only the observable payload, return codes, and latency spans.

% ==============================================================================
\vspace{12pt}
\needspace{8\baselineskip}
\section{Observability and Tracing}

\subsection{Tracing as Forensic Evidence, Not Dashboard Decoration}
In distributed AI workflows, debugging requires forensic reconstruction of the entire execution tree:
\begin{center}
\resizebox{0.92\textwidth}{!}{%
\begin{tikzpicture}[
  every node/.style={font=\small\ttfamily},
  box/.style={draw=secondary, fill=white, rounded corners=2pt, align=left, inner sep=4pt}
]
\node[box] (root) at (0, 0) {Episode: run-4a15a9 (invoke\_agent) [status=COMPLETED]};
\node[box] (c1) at (0.8, -0.8) {$\llcorner$ Step 1: chat (deepseek-v4-flash) [24ms, in=586, out=42]};
\node[box] (t1) at (1.6, -1.6) {$\llcorner$ Tool: search("tok-89f4b1") [status=OK, docs=2]};
\node[box] (c2) at (0.8, -2.4) {$\llcorner$ Step 2: chat (deepseek-v4-flash) [52ms, in=1240, out=38]};
\node[box] (t2) at (1.6, -3.2) {$\llcorner$ Tool: lookup("link-7a8f09") [status=TIMEOUT (F2), dur=120ms]};
\node[box, draw=accent] (r1) at (2.4, -4.0) {$\llcorner$ Recovery: bounded\_retry [attempt=1, backoff=10ms]};
\node[box] (t2_ret) at (2.4, -4.8) {$\llcorner$ Tool: lookup("link-7a8f09") [status=OK, dur=18ms]};
\node[box] (c3) at (0.8, -5.6) {$\llcorner$ Step 3: chat (terminal answer) [status=PASS, cited=2]};
\end{tikzpicture}%
}
\end{center}

\subsection{Crash-Surviving SQLite Spans}
Standard in-memory logging libraries lose telemetry when a fatal unhandled exception or OOM event crashes the runtime. FAULTLINE persists spans synchronously to SQLite (\code{trace.db}) with foreign-key linked parent-child spans. When a model call throws a fatal error, the terminating span records \code{OutcomeStatus.MODEL\_FAILURE} before the exception unwinds.

\subsection{OpenTelemetry GenAI Semantic Conventions (P02)}
To prevent proprietary lock-in, Project \textbf{P02} maps internal SQLite traces into standard \textbf{OpenTelemetry GenAI Semantic Conventions v1.29.0}:
\begin{itemize}
    \item Root span: \code{gen\_ai.operation.name = "invoke\_agent"}, \code{gen\_ai.system = "faultline"}.
    \item Model spans: \code{gen\_ai.operation.name = "chat"}, capturing \code{gen\_ai.request.model}, \code{gen\_ai.usage.input\_tokens}, and \code{gen\_ai.usage.output\_tokens}.
    \item Tool spans: \code{gen\_ai.operation.name = "execute\_tool"}, capturing \code{gen\_ai.tool.name}.
\end{itemize}
\textbf{Conformance Verification:} \code{tests/phase2/test\_otel\_conformance.py} verifies that all 5,392 spans across 350 exported traces achieve \textbf{100.0\% conformance} against the official SemConv v1.29.0 JSON schema with \textbf{zero attribute drift}.

\subsection{Google SRE Multi-Window Multi-Burn-Rate Alerting (P13)}
Project \textbf{P13} integrates Google SRE-style multi-burn-rate alerting over the \code{trace.db} store:
\begin{itemize}
    \item \textbf{SLO Definition:} $95.0\%$ Grounded Pass Rate over a 30-day compliance window ($720\text{ hours}$). Error budget = $5.0\%$.
    \item \textbf{Fast-Burn Alert (Page):} Triggers when the 1-hour burn rate exceeds $14.4\times$ (consuming $2.08\%$ of monthly budget in 1 hour).
    \item \textbf{Empirical Drill Result:} During an injected tool protocol failure (\code{INC-20260904-01}), pass rate collapsed from $100\%$ to $25\%$ ($75\%$ error rate). The evaluator detected a \textbf{$15.0\times$ fast-burn rate} within 6 minutes, firing a page alert and documenting recovery after rollback.
\end{itemize}

% ==============================================================================
\vspace{12pt}
\needspace{8\baselineskip}
\section{Failure Detection}

\subsection{Detection vs Diagnosis vs Correctness}
Reliability engineering requires separating four sequential questions:
\begin{enumerate}
    \item \textbf{Detection:} Did an anomaly or failure occur? (Binary signal).
    \item \textbf{Diagnosis:} Which component failed and what was the fault class? (F1--F6 attribution).
    \item \textbf{Correctness Evaluation:} Is the resulting answer accurate and grounded? (Oracle gate).
    \item \textbf{Recovery Evaluation:} Did the recovery action restore user-visible correctness within budget?
\end{enumerate}

\subsection{The Q2 Detection Findings: Aggregate Lies}
In Day 15 (Q2), the deterministic detector was evaluated across 44 labelled test runs:
\begin{center}
\small
\begin{tabularx}{\textwidth}{p{0.36\textwidth}cccc}
\toprule
\textbf{Detector Evaluation Slice} & \textbf{Precision} & \textbf{Recall} & \textbf{F1 Score} & \textbf{Miss Breakdown} \\
\midrule
Overall Aggregate & \textbf{1.000} & \textbf{0.615} & \textbf{0.762} & 10 False Negatives \\
Structural Transport (F1, F2, F4, F6) & 1.000 & 1.000 & 1.000 & 0 False Negatives \\
Semantic Escapes (F3, F5) & 1.000 & 0.200 & 0.333 & 10 False Negatives \\
\bottomrule
\end{tabularx}
\end{center}

\begin{callout}[accent]{The Anatomy of the 10 False Negatives}
Of the 10 false negatives surfaced in the Q2 audit:
\begin{itemize}
    \item \textbf{6 misses were threshold-reducible:} Latency or repetition spikes that sat slightly below hardcoded heuristics.
    \item \textbf{4 misses were semantic escapes:} Valid tool outputs containing plausible but factually corrupted context that zero deterministic schema checks could detect.
\end{itemize}
This finding proves that \textbf{aggregate F1 scores conceal dangerous operational blind spots}.
\end{callout}
